\documentclass[11pt]{article}
\usepackage{amsmath,amssymb,amsthm}
\newtheorem{proposition}{Proposition}
\title{A Small Exact Test of Reverse Sign Rank}
\author{}
\date{November 23, 2026}
\begin{document}
\maketitle

\paragraph{Scope.}
The conjectured bound $\operatorname{signrank}(\operatorname{sign}F)\leq2H$
for a reverse-admissible $H$-head grid score is false. This note supplies an
explicit rational certificate using two heads. It neither disproves an
unspecified $O(H)$ bound nor proves such a bound. In particular, it is not
an obstruction to the requested logarithmic saving. All results here are
written proofs and exact rational tests, not Lean declarations.

\begin{proposition}
There is a two-head score on the Boolean $37$-cube with strictly negative
denominator slopes and denominator minima greater than $9/10$, whose
$5$ by $32$ grid sign matrix has sign rank exactly $5$.
\end{proposition}
\begin{proof}
Use the five row coefficients in
\texttt{reverse-joint-five-row.tex}, with $s=32$ and
$\varepsilon=1/3300$. Write its two poles as
$d_i=1-c_i\varepsilon$ and $e_i=1-k_i\varepsilon$, and its two numerator
coefficients as $A_i,C_i$. Thus
\[
 G_i(u,v,a,b)=(A_i+a)(e_i-v)+(C_i+b)(d_i-u)
\]
is the cleared score. The two rational points
$x_+=(\varepsilon,2\varepsilon,3,5)$ and
$x_-=(2\varepsilon,\varepsilon,5,3)$ satisfy $G_i(x_\tau)=0$ for
$1\leq i\leq4$ and $\tau G_5(x_\tau)>0$.
The displayed determinant formula in that note proves that each first-four
Jacobian $J_\tau$ is invertible over the rationals. No inverse function
theorem is needed for the following construction.

Given $\sigma\in\{-1,1\}^4$, put
$v=J_\tau^{-1}\sigma$ and $q=-v_a v_v-v_b v_u$.
For every row the quadratic remainder is the same, so
\[
 G_i(x_\tau+t v)=t\sigma_i+t^2q\quad(1\leq i\leq4),
 \qquad
 G_5(x_\tau+t v)=g+t\ell+t^2q,
\]
where $g=G_5(x_\tau)$ and
$\ell=\nabla G_5(x_\tau)\cdot v$ are rational.
Choose the following explicit positive rational step:
\[
 t=\min\left\{
 \frac12,
 \frac{\varepsilon}{4(1+|v_u|+|v_v|)},
 \frac1{4(1+|q|)},
 \frac{|g|}{4(1+|\ell|+|q|)}
 \right\}.
\]
Then $|tq|<1/4$, so the first four signs are $\sigma$.
Also $t\leq1$ and
$|t\ell+t^2q|\leq t(|\ell|+|q|)<|g|/4$, so the fifth sign is
$\tau$. The denominator parameters remain in
$\varepsilon/2<u,v<5\varepsilon/2$.

Apply this construction once for each of the $32$ pairs $(\sigma,\tau)$,
using these rational points as the columns. The same full-cube denominator
calculation as in the five-row construction gives, for either head,
\[
 \min B>1-37\varepsilon-32\frac52\varepsilon
 =1-\frac{117}{3300}>\frac9{10}.
\]
All row and column slopes are strictly negative. Positive denominator
products preserve the signs of the cleared scores, so the resulting
matrix contains every column of $\{-1,1\}^5$ exactly once.

Every real matrix with these strict signs has linearly independent rows.
Indeed, if its rows obeyed a nonzero relation with coefficients $a_i$,
choose a column whose $i$th sign agrees with $a_i$ whenever $a_i\neq0$.
The corresponding entry of the alleged relation is a sum of positive
numbers, a contradiction. Thus its rank is at least $5$. There are only
five rows, proving that the sign rank is exactly $5$.
\end{proof}

\paragraph{Verification.}
\texttt{reverse-small-signrank-check.py} constructs all $32$ rational
columns by the displayed formula, verifies the exact quadratic expansion,
checks all $160$ strict signs, and checks both whole-cube minima. It uses
neither floating-point arithmetic nor a numerical solver. The existing
regular-sheet construction with path-comparison readouts gives the stronger
asymptotic ratio $(3H-1)/(H+1)$, but that ratio still remains bounded.

\paragraph{What the approximation argument does not establish.}
The reverse coefficient budget allows rapidly convergent denominator
expansions away from a small set of large coefficients. Truncating such
expansions gives low-rank approximations. There is no uniform positive
lower bound on the strict margins in this model, however, so an error
bound alone does not preserve all signs. Neither low displacement rank
of individual Cauchy or Loewner matrices nor ordinary parameter-side VC
dimension resolves this issue. This audit establishes no universal
linear sign-rank bound and no superlinear reverse sign-rank family.

\end{document}
